% Owner: classification_e.  Independent draft, not yet input during draft19 freeze.
% Referees requested: root and boundary_d.  No arbitrary-chain conclusion.
\begin{proposition}[E.2.7 (Proved): an extended fibration cannot reach the zero period]
Under the hypotheses of E.2.5, suppose in addition that the
intrinsic algebraic reductions extend to a holomorphic map
\[
 g:\mathcal X\longrightarrow\mathbb P^1\times\Delta
\]
whose restriction to $X_0$ is nonconstant.  This hypothesis
holds in E.2.6.  Normalize the punctured family to fixed
discrete logarithmic data with multiplicities $(3k,4l)$, and
write its isomorphism parameter as
\[
 z=\exp\!\left(-\frac{\pi i c}{kl}\right).
\]
The resulting holomorphic function $z:\Delta^*\to\mathbb C^*$
extends across zero and satisfies $z(0)\ne0$.  Its central
value lies strictly inside the admissible period disc.

This statement concerns the period function and the extended
relative fibration.  It does not identify $X_0$ with the
member at $z(0)$, and does not assert separatedness of
different smooth completions of the same punctured family.
\end{proposition}
\begin{proof}
Local miniversality in E.2.1 supplies holomorphic local
classifying functions for the family on $\Delta^*$.
E.1.2--E.1.3 normalize the discrete data by holomorphic
family comparisons.  For fixed data the comparison shifts
are precisely $c\mapsto c+2kl\mathbb Z$, so the displayed
exponential is a single-valued holomorphic function on
$\Delta^*$.  An automorphism with the same marked parameter
has identity linear action on the homology of an ordinary
torus: in the necessity argument of E.1.2 its parameters
are $\alpha=1$ and $h=0$.  Consequently the homology local
system of the punctured family is also the one obtained
from these local classifying maps; possible translation
automorphisms of the comparisons add no homology monodromy.

The admissible parameter domain is a punctured disc
$0<|z|<R_{k,l}$, with finite positive radius.  The removable
singularity theorem therefore extends $z$ holomorphically
to $\Delta$.  It remains to exclude a zero at the centre.

We first control the moving three critical values without
assuming that they remain distinct at zero.  Since
$\pi:\mathcal X\to\Delta$ is a submersion, the critical
locus of $g$ is the analytic locus where its differential
has rank less than two.  Properness makes its image
\(\mathcal C\subset\mathbb P^1\times\Delta\) a closed
analytic set.  Its intersection with the central target
curve is the finite critical-value set of the nonconstant
map $g_0:X_0\to\mathbb P^1$.  In particular
$\mathcal C$ has no component equal to that whole curve.
After shrinking $\Delta$, its reduced one-dimensional
part is finite over $\Delta$ and over $\Delta^*$ consists
of the three critical-value sections.

Those sections are individually marked by the linear
peripheral monodromy orders $3$, $4$, and infinity.
The normalization of the closure of each section is
therefore finite of degree one over the disc.  It is
isomorphic to the disc, so the three values
$c_1(s),c_2(s),c_0(s)$ extend as holomorphic maps to
$\mathbb P^1$, with collisions allowed at $s=0$.
The unique projective transformation $A_s$ sending this
ordered triple to $(0,1,\infty)$ has meromorphic
coefficients at zero, by its explicit cross-ratio formula.
This is the critical-value argument of B.4.4 in a setting
where no vertical discriminant component is present.

Choose a point $v\in\mathbb P^1$ outside all central
critical values.  Properness and closedness of
$\mathcal C$ allow us to shrink the disc so that
$(v,s)\notin\mathcal C$ for every $s$.  Then
\[
 \mathcal T=g^{-1}(\{v\}\times\Delta)\longrightarrow\Delta
\]
is a proper smooth holomorphic family of surfaces.
Its noncentral fibres are ordinary complex two-tori, and
Ehresmann's theorem makes its integral first-homology
monodromy around $s=0$ the identity.

In the original marked base, this same family is obtained
by evaluating the period construction at
\[
 t(s)=A_s(v)\in\mathbb P^1\setminus\{0,1,\infty\}.
\]
The function $t(s)$ is meromorphic at zero.  Its loop on a
sufficiently small positive $s$-circle is therefore either
null-homotopic in the marked base or conjugate to a power
of one of its three peripheral loops.  Indeed, if its
limit is ordinary, its image lies in an ordinary small
disc; if its limit is a marked point, a local coordinate
has a zero or pole of finite order and its image lies
in the corresponding punctured disc.  This observation
also covers a constant $t(s)$.

Let $\Lambda$ have ordered basis
$(\widehat\gamma,\widehat u,\widehat w,\widehat\delta)$,
put $E=\widehat\delta\otimes\gamma$, and write
$S_a=I+aE$.  The two finite peripheral monodromies are
\[
 A_1=\begin{pmatrix}
 1&0&0&0\\6&0&1&0\\-6&-1&-1&0\\-2&1&0&1
 \end{pmatrix},\qquad
 A_2=\begin{pmatrix}
 1&0&0&0\\0&0&-1&0\\-6&1&0&0\\3&0&1&1
 \end{pmatrix}.
\]
They have orders three and four, respectively, even for
the larger multiplicities: the extra finite deck
transformations are translations and do not affect
$H_1$.  In the cusp marking the third monodromy is
\[
 A_0=I+N,\qquad
 N\widehat u=\widehat w,\quad
 N\widehat\gamma=-\widehat\delta,\quad
 N\widehat w=N\widehat\delta=0.
\]
All peripheral matrices preserve $\widehat\delta$ and
$\gamma$, so commute with every $S_a$.  The change to
the cusp marking preserves the same rank-one operator
$E$; conjugating a peripheral loop consequently does
not alter the cancellation test below.

Suppose that $z(s)$ has a zero of order $d>0$ at zero.
Continuation of a logarithm around the positive loop
changes $c$ by $-h$, where
\[
 h=2kl\,d>0.
\]
Since $\Pi_{c+h}=\Pi_cS_h$, a comparison of the identical
complex universal covers from parameter $c-h$ back to
$c$ acts on the marked lattice by $S_{-h}$.
The affine translation part of a closing comparison
acts trivially on $H_1$.  Thus the return map of the
family $\mathcal T$ has homology action, up to a common
conjugation,
\[
 S_{-h},\qquad S_{-h}A_1^n,\qquad
 S_{-h}A_2^n,\qquad\hbox{or}\qquad S_{-h}A_0^n
\]
for an integer $n$.  Reversing the convention for a
peripheral generator only replaces $n$ by $-n$.
The parameter factor and the base factor commute, so
the order in this display has no effect.

None of these matrices is the identity.  In the two
finite cases, taking the third or fourth power would
give $S_{-3h}=I$ or $S_{-4h}=I$, impossible because
$E\ne0$ and $h>0$.  In the cusp case, $N^2=EN=NE=0$
and hence
\[
 S_{-h}A_0^n=I+nN-hE.
\]
Its coefficient of $\widehat w$ on $\widehat u$ is
$n$, so identity would force $n=0$, and then $h=0$.
The ordinary case is $S_{-h}\ne I$.  This contradicts
the identity monodromy of the smooth family over
the full disc.  Therefore $z(0)\ne0$.

Finally, the extended holomorphic function satisfies
$|z(0)|\le R_{k,l}$.  Equality would make it a constant
of maximal modulus, contradicting $|z(s)|<R_{k,l}$
for $s\ne0$.  Thus $0<|z(0)|<R_{k,l}$, as claimed.
The matrix identities and the factor $h=2kld$ are
checked exactly in \path{code/e_push_E_zero_end.py}.
\end{proof}

\begin{remark}[Remaining scope after E.2.7]
For positive canonical slope, E.2.6--E.2.7 reduce a
one-sided smooth specialization to an interior period
value.  Equality of that value with an existing model
does not by itself prove uniqueness of its smooth
completion.  For negative slope the horizontal fixed
divisor in E.2.5 still prevents applying E.2.7 without
an additional fibration-persistence argument.  Neither
restriction addresses an arbitrary new deformation
branch through an unconstructed smooth limit.
\end{remark}
